===== KEY: v2:1e804bc1faa54fb620091e3187ccc58a1eef5904bb0e702085fa8c6939e0f14d =====
% ============================================================
% FRAGMENT 1
% ============================================================
\subsection{Low-energy structure of $W^{00}$: axial transitions and a parameter-free test}
\label{subsec:W00lowE}

The wavepacket hadronic tensor $W^{00}$ exhibits a pronounced ridge at frequencies far below the two-meson threshold $\omega=5.490$. Its origin is fixed by symmetry. The lightest meson is a $\rho^{0}$ analog: charge conjugation forces the elastic vector form factor to vanish identically, $\langle k'|J^{0}_{V}|k\rangle=0$ within the band \cite{placeholder}, so no intra-band vector response survives. The ridge is instead built from the axial channel, whose forward charge is large, $\langle k|J_{\mathrm{ax}}(0)|k\rangle=7.24\text{--}7.55$ across the band. Two mechanisms contribute: intra-band beats at $\omega_{k}(q_{1})=E(k{+}q_{1})-E(k)$, generated by the coherence of the packet, and interband $M\!\to\!M'$ transitions into the second meson band at gaps $3.155$--$3.285$. Both are captured by a two-band correlator whose inputs---packet amplitudes $c_{k}$, band energies, and axial matrix elements---are all determined independently, leaving no fit parameters:
\begin{equation}
W^{00}_{\mathrm{2b}}(q_{1},t)=\sum_{k}\Big[\,c^{*}_{k+q_{1}}c_{k}\,
\mathcal{A}_{k}(q_{1})\,e^{-i[E(k+q_{1})-E(k)]t}
+|c_{k}|^{2}\,\mathcal{B}_{k}(q_{1})\,e^{-i[E'(k+q_{1})-E(k)]t}\Big],
\label{eq:twoband}
\end{equation}
with $\mathcal{A}$ ($\mathcal{B}$) the intra-band (interband) axial-current matrix elements and $E'$ the second-band dispersion.

Because a wavepacket is not a Hamiltonian eigenstate, the hermiticity completion $W(-t)=W^{*}(t)$ commonly used for vacuum correlators is invalid here: it scatters higher-sector spectral content into low $\omega$. We therefore employ one-sided Gaussian-windowed transforms ($\sigma_{t}=8/3$) of the real-time data, with the realized packet width $\sigma_{x}=0.83$, slightly broadened from the $\sigma_{x}=0.75$ Gaussian target. Figure~\ref{fig:ridge} overlays Eq.~\eqref{eq:twoband} on the 101-qubit data for eight curves---rest and boosted ($\bar{k}=2\pi/5$) packets at $q_{1}=\pm2\pi/5,\pm4\pi/5$---with root-mean-square deviations of $0.7$--$2.2\%$. The model also reproduces the kinematic asymmetry of the boosted packet, whose $\pm q_{1}$ peak splitting is $0.07$ versus $0.45$ at the two momentum transfers, and its measured group velocity $+0.093$ agrees with the dispersion-slope prediction $0.09$--$0.12$. An integrated (Simpson-rule) Ward-identity check on the rest-frame data closes at $6.6\%$, consistent with the $\delta t=0.5$ time-integration error that dominates it.

One systematic requires care: the sequential gate sweep used in state preparation is chiral and leaves a ${\sim}13\%$ reflection asymmetry in the raw rest-frame correlator. We mitigate it by forming observables odd in $q_{1}$, for which the rest packet furnishes a null test, and the quoted agreement is quoted after this construction \cite{placeholder}.

% ============================================================
% FRAGMENT 2
% ============================================================
\subsection{Elastic meson--meson phase shifts from the same dataset framework}
\label{subsec:phaseshifts}

The same variational-plus-real-time framework yields two-meson finite-volume energies above the noninteracting threshold $E_{\mathrm{thr}}=5.490$, from which elastic phase shifts follow by the $(1{+}1)$-dimensional analog of the L\"uscher quantization condition \cite{placeholder}. For two identical mesons with back-to-back momenta $\pm p$ on a periodic volume of length $L$, single-valuedness of the wavefunction requires
\begin{equation}
p_{n}L+2\delta(p_{n})=2\pi n,\qquad E_{n}=E(p_{n})+E(-p_{n}),
\label{eq:luescher}
\end{equation}
where $E(k)$ is the single-meson dispersion measured in Sec.~\ref{subsec:W00lowE}, so that the shift of $E_{n}$ from its noninteracting value directly determines $\delta(p_{n})$. Applying Eq.~\eqref{eq:luescher} to the measured levels gives
$\delta(p)=+0.56,\;+0.15,\;+0.67$ at $p=1.03,\;1.50,\;2.81$, uniformly positive and hence attractive.

The $n_{s}=6$ level is excluded from this analysis: it sits at ${\approx}5.00$, degenerate with the two flat stable bands at $M^{*}\simeq5.00$, identifying it as a stable single-particle excitation rather than a two-meson scattering state. Equation~\eqref{eq:luescher} presumes elastic two-body kinematics and would assign this level a spurious phase shift; its flat (volume-independent) dispersion confirms the one-body assignment. The remaining levels lie safely between $E_{\mathrm{thr}}$ and the onset of deeper inelasticities.

These phase shifts, obtained without additional circuit executions, convert the quantum dataset into falsifiable output: Eq.~\eqref{eq:luescher} generates prediction tables for two-meson spectra at other volumes, providing sharp targets for future hardware runs \cite{placeholder}.
